\documentclass[11pt]{article}

\usepackage[T1]{fontenc}
\usepackage{lmodern}
\usepackage{microtype}
\usepackage[a4paper,margin=29mm]{geometry}
\usepackage{amsmath,amssymb,amsthm,mathtools}
\usepackage{booktabs}
\usepackage{enumitem}
\usepackage{tikz}
\usetikzlibrary{arrows.meta,positioning,shapes.geometric}
\usepackage[colorlinks=true,allcolors=blue!55!black]{hyperref}
\hypersetup{
  pdftitle={The third magnetic Neumann eigenvalue does not exceed the Dirichlet ground state on an annulus},
  pdfauthor={Xiaojian Zeng},
  pdfkeywords={magnetic Laplacian, annulus, Neumann-Dirichlet comparison, homogeneous magnetic field, min-max principle}
}

\newtheorem{theorem}{Theorem}[section]
\newtheorem{proposition}[theorem]{Proposition}
\newtheorem{lemma}[theorem]{Lemma}
\newtheorem{remark}[theorem]{Remark}

\newcommand{\Aa}{\mathcal A_a}
\newcommand{\Bdisk}{\mathbb B}
\newcommand{\R}{\mathbb R}
\newcommand{\Z}{\mathbb Z}
\newcommand{\N}{\mathbb N}
\newcommand{\dd}{\,\mathrm d}
\newcommand{\Afield}{\mathbf A_b}
\newcommand{\qN}{q^{N}_{a,b}}
\newcommand{\qD}{q^{D}_{a,b}}

\title{The third magnetic Neumann eigenvalue does not exceed the Dirichlet ground state on an annulus}
\author{Xiaojian Zeng\\
\small PBC School of Finance, Tsinghua University, Beijing 100190, China\\
\small \href{mailto:xiaojianzeng@foxmail.com}{xiaojianzeng@foxmail.com}}
\date{}

\begin{document}
\maketitle

\begin{abstract}
Let $\Aa=\{x\in\R^2:a<|x|<1\}$, $0<a<1$, and consider the
two-dimensional magnetic Laplacian $(-i\nabla-\Afield)^2$ generated by the
homogeneous field in symmetric gauge, $\Afield=b(-x_2,x_1)/2$, $b>0$.
Both boundary circles carry either magnetic Neumann or Dirichlet conditions,
and no independent Aharonov--Bohm flux is added.  We prove that the third
Neumann eigenvalue, counted with multiplicity, never exceeds the Dirichlet
ground-state energy.  Thus the additional boundary component and the
associated rearrangement of angular momenta do not destroy the
third-Neumann/first-Dirichlet ordering known numerically for the disk.
The proof is exact and partitions the field--radius plane into three regimes.
For small fields, a quotient-monotonicity argument transfers three disk-type
trials to every annulus.  Intermediate fields are controlled by restricting
two explicit disk modes at the exactly solvable field $b=2$.  At high field,
Landau modes settle flux windows containing three integers; the complementary
windows are handled by moment bounds for thick annuli and transverse
Dirichlet confinement for thin annuli.  No finite angular cutoff,
perturbative truncation, or numerical eigenvalue enclosure enters the proof.
\end{abstract}

\noindent\textbf{Keywords:} magnetic Laplacian; annulus; Neumann--Dirichlet
eigenvalue comparison; homogeneous magnetic field; min--max principle

\noindent\textbf{Mathematics Subject Classification:} 35P15; 81Q10

\section{Introduction}

Comparison inequalities between Dirichlet and Neumann spectra are sensitive
to geometry, topology, and magnetic flux.  For the homogeneous magnetic
Laplacian on a bounded convex planar domain, Lotoreichik proved
$\mu_{k+1}\leq\lambda_k$ for every $k$ \cite{Lotoreichik2025}.  The same
paper analyzed the unit disk by angular fibers: at the first Dirichlet
threshold it obtained an analytic comparison for $b>4$, numerical evidence
for the remaining fields, and an exact crossing at $b=2$.  Earlier work of
Frank and Laptev gives a one-level counting advantage at Landau thresholds
on general domains \cite{FrankLaptev2010}.  These results do not directly
give a two-level shift at an arbitrary Dirichlet threshold on an annulus,
which is nonconvex and has two boundary components.

The topology also forces a convention that is immaterial on simply
connected domains: adding a closed, non-exact one-form changes the flux and
can change the spectrum \cite{ColboisEtAl2023}.  Throughout this article we
fix the standard symmetric potential for the homogeneous field and add no
independent harmonic Aharonov--Bohm component.  General background on the
constant-field magnetic Laplacian can be found in
\cite{FournaisHelffer2010}.

Our result establishes the following annular comparison in this fixed model.

\begin{theorem}\label{thm:main}
For every $0<a<1$ and $b>0$, the globally ordered eigenvalues of the
homogeneous-field magnetic Laplacian on $\Aa$, counted with multiplicity,
satisfy
\[
        \mu_3(\Aa,b)\leq \lambda_1(\Aa,b),
\]
where $\mu_j$ denotes the spectrum with magnetic Neumann conditions on both
circles and $\lambda_j$ the spectrum with Dirichlet conditions on both
circles.
\end{theorem}

The proof is variational.  It constructs three mutually orthogonal angular
trials below explicit lower bounds for the Dirichlet ground state.  The
choice of trials changes at $b=2$ and $b=6$, and at high field it also depends
on the number of integers in the annular flux window.  Figure
\ref{fig:regime-map} records the exhaustive logical partition; the estimates
themselves are proved in Sections \ref{sec:lower}--\ref{sec:high}.

\input{manuscript/Fig1.tex}

In the chosen dimensionless normalization, the magnetic flux through the
annulus divided by $2\pi$ is $b(1-a^2)/2$.  The theorem is therefore a
spectral statement about an idealized one-particle model.  It is not an
assertion about a measured mesoscopic ring, a many-body system, or an added
Aharonov--Bohm flux.

\section{Operator, fibers, and a min--max criterion}\label{sec:model}

Let
\[
 \Aa=\{(r\cos\theta,r\sin\theta):a<r<1\},\qquad
 \Afield(x)=\frac b2(-x_2,x_1),
\]
with $0<a<1$ and $b>0$.  The Neumann and Dirichlet realizations of
\[
 H_{a,b}=(-i\nabla-\Afield)^2
\]
are associated with the same quadratic expression
\[
 q_{a,b}[u]=\int_{\Aa}|(-i\nabla-\Afield)u|^2\dd x
\]
on $H^1(\Aa)$ and $H^1_0(\Aa)$, respectively.  Because $\Afield$ is tangent
to both boundary circles, the magnetic Neumann condition is the ordinary
radial Neumann condition in every angular fiber.

For $u(r,\theta)=f(r)e^{im\theta}$, $m\in\Z$, the radial expression in
$L^2((a,1);r\dd r)$ is
\begin{equation}\label{eq:fiber}
 L_{m,b}f=-f''-\frac1r f'
       +\left(\frac mr-\frac{br}{2}\right)^2f.
\end{equation}
The global spectra are the ordered union of these fibers, with degeneracies
retained.

\begin{lemma}[Three-trial criterion]\label{lem:minmax}
Suppose $u_1,u_2,u_3\in H^1(\Aa)$ have distinct angular indices and each has
Rayleigh quotient at most $L$.  Then $\mu_3(\Aa,b)\leq L$.  In particular,
if $L\leq\lambda_1(\Aa,b)$, the conclusion of Theorem \ref{thm:main}
follows.
\end{lemma}

\begin{proof}
Distinct angular sectors are orthogonal both in $L^2$ and for the magnetic
quadratic form.  Hence every nonzero vector in their three-dimensional span
has quotient bounded by the largest of the three individual quotients.  The
min--max principle gives the claim.
\end{proof}

\section{Three lower bounds for the Dirichlet ground state}\label{sec:lower}

We use complementary bounds adapted to low field, Landau windows, and thin
rings.

\begin{proposition}\label{prop:lower}
For all $0<a<1$ and $b>0$,
\begin{align}
 \lambda_1(\Aa,b)&\geq
 \begin{cases}
  5+b^2/4,&0<b\leq2,\\
  6,&b\geq2,
 \end{cases}                                                     \label{eq:barta}\tag{2.1}\\
 \lambda_1(\Aa,b)&\geq b,                                      \label{eq:landau-lower}\tag{2.2}\\
 \lambda_1(\Aa,b)&\geq W(a):=
   \frac{\pi^2}{(1-a)^2}-\frac1{4a^2}.                         \label{eq:width}\tag{2.3}
\end{align}
\end{proposition}

\begin{proof}
Let $\Bdisk$ be the unit disk.  Dirichlet domain monotonicity gives
$\lambda_1(\Aa,b)\geq\lambda_1(\Bdisk,b)$.  The Dirichlet disk ground state
is radial for the homogeneous field \cite[Section 5]{Lotoreichik2025}.
For the positive radial comparison function
\[
       \phi(r)=(1-r^2)e^{-r^2/2}
\]
one computes
\[
       \frac{L_{0,b}\phi}{\phi}
       =6+\frac{b^2-4}{4}r^2.
\]
Barta's inequality, followed by minimization over $0<r<1$, proves
\eqref{eq:barta}.

For $f(a)=f(1)=0$, direct completion of the square in any $m$-fiber gives
\begin{equation}\label{eq:factorization}
 q^D_{m,a,b}[f]-b\|f\|^2
 =2\pi\int_a^1 r\left|f'-
       \left(\frac mr-\frac{br}{2}\right)f\right|^2\dd r\geq0.
\end{equation}
This proves \eqref{eq:landau-lower} after summing the angular fibers.

Finally set $g=r^{1/2}f$.  The endpoint conditions give
\[
 \int_a^1r|f'|^2\dd r
 =\int_a^1\left(|g'|^2-\frac{|g|^2}{4r^2}\right)\dd r
 \geq\left(\frac{\pi^2}{(1-a)^2}-\frac1{4a^2}\right)
       \int_a^1|g|^2\dd r.
\]
Dropping the nonnegative magnetic potential proves \eqref{eq:width}.
\end{proof}

\section[Small fields: 0 < b <= 2]{Small fields: $0<b\leq2$}\label{sec:small}

Put $c=b/2\in(0,1]$ and $x=a^2$.  Consider the three angular trials
\begin{equation}\label{eq:small-trials}
 1,\qquad r e^{-cr^2/2}e^{i\theta},\qquad
 r e^{-cr^2/2}e^{-i\theta}.
\end{equation}
Write their quotients as $R_0,R_1,R_{-1}$.  Direct integration gives
\begin{equation}\label{eq:R0-small}
 R_0=\frac{c^2(1+x)}2<5+c^2,
 \qquad R_{-1}=R_1+4c.
\end{equation}
The $m=1$ radial factor satisfies $L_{1,2c}f=2cf$.  If
\[
 I(x,c)=\int_a^1r^3e^{-cr^2}\dd r
 =\frac{(1+cx)e^{-cx}-(1+c)e^{-c}}{2c^2},
\]
integration by parts yields
\begin{equation}\label{eq:R1-exact}
 R_1(x,c)=2c+
 \frac{(1-c)e^{-c}-x(1-cx)e^{-cx}}{I(x,c)}.
\end{equation}
At the disk endpoint $x=0$ this implies
\begin{equation}\label{eq:disk-gap}
 5+c^2-R_{-1}(0,c)
 =(1-c)\left(5-c-\frac{2c^2}{e^c-1-c}\right).
\end{equation}
For $0<c\leq1$, positivity of the first four relevant terms in the
exponential series gives
\begin{equation}\label{eq:disk-gap-positive}
 (5-c)(e^c-1-c)-2c^2
 \geq c^2\left(\frac12+\frac c3-\frac{c^2}{6}\right)>0.
\end{equation}
Thus the disk gap in \eqref{eq:disk-gap} is positive when $c<1$.

It remains to show that removing a centered disk cannot increase $R_1$.
Differentiation of \eqref{eq:R1-exact} gives a positive prefactor times
$P(x,c)$:
\begin{equation}\label{eq:derivative-P}
 \partial_xR_1=
 \frac{2c^2e^c}{(cxe^c-ce^{cx}+e^c-e^{cx})^2}\,P(x,c),
\end{equation}
where
\begin{align*}
P(x,c)={}&c^3x^2e^{cx}-c^3xe^{cx}+c^2x^2e^c+c^2x^2e^{cx}
 -2c^2xe^{cx}\\
&+2cxe^c-3cxe^{cx}+ce^{cx}-e^c+e^{cx}.
\end{align*}
Set $t=1-x$ and $H(t,c)=e^{ct-c}P(1-t,c)$.  Exact differentiation gives
\begin{align}
 H(0,c)&=H_t(0,c)=0,\qquad
 H_{tt}(0,c)=c^2(c^2-1)\leq0,                                  \label{eq:Hdata}\\
 H_{ttt}(t,c)&=c^3e^{ct}
 \bigl[c^2(1-t)^2-4c(1-t)-1\bigr]<0                            \label{eq:Hthird}
\end{align}
for $0<t\leq1$ and $0<c\leq1$.  Three integrations from $t=0$ show
$H(t,c)<0$, hence $\partial_xR_1<0$ for $0\leq x<1$.

For $c<1$, equations \eqref{eq:disk-gap}--\eqref{eq:Hthird} put both signed
angular quotients strictly below $5+c^2$.  At $c=1$ the disk gap vanishes,
but the strict decrease for every annulus $a>0$ gives the same conclusion.
Together with \eqref{eq:R0-small}, Proposition \ref{prop:lower}, and Lemma
\ref{lem:minmax}, this proves Theorem \ref{thm:main} for $0<b\leq2$.

\section[Intermediate fields: 2 <= b <= 6]{Intermediate fields: $2\leq b\leq6$}\label{sec:intermediate}

Use the constant and the two angular trials
\begin{equation}\label{eq:mid-trials}
 f_1(r)e^{i\theta},\quad f_1(r)=re^{-r^2/2},
 \qquad
 f_2(r)e^{2i\theta},\quad f_2(r)=r^2(3-r^2)e^{-r^2/2}.
\end{equation}
At $b=2$ these are exact disk eigenfunctions:
\[
 L_{1,2}f_1=2f_1,\qquad L_{2,2}f_2=6f_2,\qquad
 f_1'(1)=f_2'(1)=0.
\]
Moreover $f_j(a)f_j'(a)>0$ for $0<a<1$.  Restriction to the annulus and
integration by parts therefore give
\begin{equation}\label{eq:restriction}
       R_1(2)<2,\qquad R_2(2)<6.
\end{equation}
Define the conditional moments
\[
 C_j(a)=\frac{\int_a^1r^3f_j(r)^2\dd r}
              {\int_a^1r f_j(r)^2\dd r},\qquad 0<C_j(a)<1.
\]
Expanding the magnetic potential around $b=2$ gives the exact shift
\begin{equation}\label{eq:field-shift}
 R_j(b)=R_j(2)+(b-2)\left[-j+\frac{b+2}{4}C_j(a)\right].
\end{equation}
Consequently,
\begin{equation}\label{eq:mid-bounds}
 R_1(b)<2+\frac{(b-2)^2}{4}\leq6,\qquad
 R_2(b)<6+\frac{(b-2)(b-6)}4\leq6.
\end{equation}
The constant quotient is
\[
       R_0=\frac{b^2(1+a^2)}8.
\]
If $a^2\leq1/3$, then $R_0\leq6$.  If $a^2>1/3$, then $a>1/2$,
$R_0<9$, and \eqref{eq:width}, with $\pi>3$, gives $W(a)>35$.  In either
case the three quotients are bounded by the Dirichlet ground-state lower
bounds.  Lemma \ref{lem:minmax} proves the theorem on this field interval.

\section[High fields: b > 6]{High fields: $b>6$}\label{sec:high}

The high-field proof first uses exact Landau trials and then treats their
integer-count complement.

\subsection{The annular flux window}

For $m\in\Z_{\geq0}$ set
\[
 u_m(r,\theta)=r^me^{-br^2/4}e^{im\theta},\qquad
 I_{a,b,m}=\int_a^1r^{2m+1}e^{-br^2/2}\dd r.
\]
The radial factor obeys
\[
 f_m'=\left(\frac mr-\frac{br}{2}\right)f_m,
 \qquad L_{m,b}f_m=bf_m.
\]
Its exact Rayleigh quotient is
\begin{equation}\label{eq:landau-trial}
 R_m=b+\frac{(m-b/2)e^{-b/2}
 -(m-ba^2/2)a^{2m}e^{-ba^2/2}}{I_{a,b,m}}.
\end{equation}
Thus $R_m<b$ whenever
\begin{equation}\label{eq:flux-window}
       \frac{ba^2}{2}\leq m\leq\frac b2.
\end{equation}
If the closed interval in \eqref{eq:flux-window} contains at least three
integers, three orthogonal trials lie below $b\leq\lambda_1$ by
\eqref{eq:landau-lower}, and the proof is complete.

Suppose that the interval contains fewer than three integers.  Every closed
interval of length at least three contains at least three integers, so
\begin{equation}\label{eq:short-window}
       b(1-a^2)<6.
\end{equation}

\subsection[The complement for a <= 1/4]{The complement for $a\leq1/4$}

The moments $C_j(a)$ from Section \ref{sec:intermediate} increase with $a$:
differentiating a truncated weighted mean shows that $C_j'(a)$ has the sign
of $C_j(a)-a^2>0$.  At $a=1/4$, the substitution $z=r^2$ gives the exact
identities
\begin{align}
 e\int_{1/16}^{1}(4/5-z)ze^{-z}\dd z
 &=\frac{17}{5}-\frac{1637}{1280}e^{15/16},                    \label{eq:moment1}\\
 e\int_{1/16}^{1}(4/5-z)z^2(3-z)^2e^{-z}\dd z
 &=\frac{344}{5}-\frac{140644469}{5242880}e^{15/16}.           \label{eq:moment2}
\end{align}
For $u=15/16$, bounding the exponential-series tail geometrically after the
quadratic term yields
\begin{equation}\label{eq:expbound}
 e^{15/16}<1+u+\frac{u^2}{2}
       +\frac{u^3/6}{1-u/4}
 =\frac{64133}{25088}<\frac{64}{25}.
\end{equation}
Substitution into \eqref{eq:moment1}--\eqref{eq:moment2} leaves, respectively,
the positive rational lower bounds
\[
       \frac{63}{500}\qquad\text{and}\qquad
       \frac{257931}{2048000}.
\]
Hence
\begin{equation}\label{eq:moment-bound}
       C_1(a)<\frac45,\qquad C_2(a)<\frac45
       \quad(0<a\leq1/4).
\end{equation}

From \eqref{eq:short-window}, $b<32/5$.  The constant trial satisfies
$b(1+a^2)<34/5<8$, hence $R_0<b$.  Equations
\eqref{eq:field-shift} and \eqref{eq:moment-bound} give
\begin{align*}
 R_1(b)&<2+\frac{17}{25}(b-2)<b,\\
 R_2(b)&<6+\frac{(b-2)(b-8)}5<6<b.
\end{align*}
All three quotients are below $b\leq\lambda_1$, which settles this subcase.

\subsection[The complement for a >= 1/4]{The complement for $a\geq1/4$}

Now use three radial-constant angular trials $e^{im\theta}$.  Put
\begin{equation}\label{eq:alpha-delta}
 \alpha=\frac{2\log(1/a)}{1-a^2},\qquad
 y=\frac{b}{2\alpha},\qquad
 \delta=\frac{1+a^2}{8}-\frac{1}{4\alpha}.
\end{equation}
Direct integration completes the square:
\begin{equation}\label{eq:constant-completion}
       R_m=\alpha(m-y)^2+b^2\delta.
\end{equation}
Choose an integer nearest to $y$ and its two neighbors.  These are three
distinct angular modes and their distances from $y$ are at most $3/2$.

The number $\alpha$ is the average of $1/z$ over $[a^2,1]$ and decreases
with $a$.  Since $e^{3/2}>4$ already follows from the cubic Taylor sum,
\begin{equation}\label{eq:alpha-bound}
       \alpha\leq\frac{32\log4}{15}<\frac{16}{5}.
\end{equation}
Also $\delta=(A-L)/4$, where $A$ and $L$ are the arithmetic and logarithmic
means of $1$ and $a^2$.  The standard mean inequalities $a\leq L\leq A$
give
\begin{equation}\label{eq:delta-bound}
       0\leq\delta\leq\frac{(1-a)^2}{8}.
\end{equation}
Using \eqref{eq:short-window} and $a\geq1/4$, the largest of the three
quotients in \eqref{eq:constant-completion} is strictly less than
\begin{equation}\label{eq:thin-upper}
 \frac{9}{2(1+a)^2}+\frac{9\alpha}{4}
 \leq\frac{72}{25}+\frac{36}{5}
 =\frac{252}{25}<12.
\end{equation}
On the other hand, \eqref{eq:width}, $a\geq1/4$, and $\pi>3$ imply
$W(a)>16-4=12$.  The three trial quotients are therefore strictly below the
Dirichlet ground state.  Together with the flux-window and small-radius
subcases, this completes the proof of Theorem \ref{thm:main}.

\section{Discussion, scope, and reproducibility}\label{sec:discussion}

The proof separates three mechanisms.  Below $b=2$, the decisive point is
monotonicity of a disk-type quotient under removal of the central region.
Between $b=2$ and $b=6$, two exactly solvable disk modes remain effective
after restriction.  Above $b=6$, either the annular flux window supplies
three below-Landau angular modes, or its shortness forces bounds that are
closed by moment control or transverse confinement.  The interfaces
$b=2$, $b=6$, and $a=1/4$ are all included in the adjacent exact estimates.

No statement is made for a nonconcentric hole, mixed boundary conditions,
an independently added harmonic flux, nonuniform fields, many-body effects,
or experimental spectra.  The inequality is non-strict in the theorem
because the proof only needs $\mu_3\leq\lambda_1$, although most variational
branches yield strict trial bounds.

All identities used above are elementary and exact.  A supplementary Python
script independently checks the Barta identity, the two explicit $b=2$
modes, the derivative factorization, the moment antiderivatives, the
constant-mode completion, and the decisive rational inequalities.  The
script is a secondary reproducibility aid; the mathematical argument does
not rely on floating-point spectra.  Regeneration instructions and the
tested software versions accompany the source bundle.

\section*{Statements and declarations}

\paragraph{Data and code availability.}
No empirical data were generated.  The exact symbolic checker and its
regeneration instructions are included with the source package.

\paragraph{Use of generative AI.}
Generative-AI tools assisted with literature triage, exploration and
adversarial checking of proof strategies, preparation of symbolic-checking
code and the regime schematic, and drafting of the manuscript.  Such tools
are not authors.  The proof is presented explicitly so that every claim can
be checked independently; responsibility for any submitted version rests
with the human author.

\paragraph{Funding and competing interests.}
No declaration is made in this local draft; funding and competing-interest
statements require confirmation by the author before submission.

\bibliographystyle{plain}
\bibliography{manuscript/references}

\end{document}
